\chapter{Literature Review}

The following chapter provides an in-depth review of foundational concepts and recent technological advances underpinning the research. The review aims to establish the theoretical base, examine pertinent frameworks and tools, and identify critical gaps justifying the present work.

\section{Theoretical Background}

\subsection{Agentic AI and Multi-Agent Systems (MAS)}

Artificial Intelligence agents are autonomous entities possessing environment perception, reasoning capabilities, and decision-making skills to pursue goals \cite{sycara1998multiagent}. Multi-Agent Systems (MAS) embody collections of such agents collaborating to solve complex, distributed tasks beyond single-agent capacity.

The Contract Net Protocol \cite{smith1980contract} remains a seminal contribution by formalizing decentralized task negotiation, allocation, and coordination among agents via bidding and task contracting, thereby laying the groundwork for distributed problem-solving architectures. The described negotiation model informs the collaborative agent interactions in the thesis, particularly in architect supervisor and domain-specific generator/validator agents.

\subsection{Retrieval-Augmented Generation (RAG)}

Large Language Models (LLMs) excel in understanding and generating human-like language but suffer from static training knowledge, causing hallucinations and inaccuracies when encountering novel or domain-specific queries.

Retrieval-Augmented Generation (RAG) \cite{lewis2020rag} mitigates such limitations by combining LLM capabilities with external knowledge sources. Upon receiving a query, the RAG framework retrieves relevant documents from curated knowledge bases, then conditions the language model's generation on the fresh, authoritative information to enhance factual correctness.

Recent enhancements like Self-RAG \cite{asai2023selfrag} introduce self-reflective mechanisms, allowing models to critique and iteratively improve their own outputs. Self-correction capabilities strongly support the iterative Evaluator-Optimizer strategy employed in the project.

\subsection{Reasoning and Acting (ReAct)}

Complex cognition requires coupling reasoning with action. The ReAct framework \cite{yao2023react} formalizes the cognition-action duality by enabling LLMs to generate explicit reasoning steps ('thoughts') interleaved with concrete actions such as tool calls or queries. The resulting cyclical reasoning-acting loop allows agents to refine results through observations and additional reasoning, forming the basis for generator and validator agents in the system.

\subsection{Multi-Agent Debate and Consensus}

Beyond single-agent reasoning frameworks like ReAct, recent work has explored whether independent language model instances can collectively reach more reliable conclusions than any single agent acting alone. Du et al.~\cite{du2023improving} demonstrate that structured multi-agent debate---where agents independently generate solutions and then critique one another over several rounds---reduces factual errors and improves reasoning quality measurably compared to single-model baselines. The intuition is that disagreement between independently initialized agents surfaces blind spots that confirmation-prone single-agent reasoning would leave unexamined. Within this framework, the compare and debate interface draws on this principle: outputs from the AWS-aligned and Azure-aligned pipelines are exposed side-by-side, allowing practitioners to reason about the relative merits of each independently derived recommendation.

\section{Implementation Frameworks}

\subsection{LangChain and LCEL}

LangChain \cite{langchain2024docs} is an extensible architecture providing building blocks for LLM applications, including prompt templates, model integration, and tool connectors. Its modern expression language, LCEL \cite{langchain2024lcel}, allows declarative specification of complex, pipelined workflows by composing modular components, enabling scalable parallel and sequential processes with ease.

\subsection{LangGraph: Enabling Agentic Loops}

While LCEL excels in linear workflows, many AI systems require iterative, stateful cycles. LangGraph \cite{langgraph2024docs} extends LangChain to support cyclic `StateGraph`s, where nodes represent agents modifying shared state, and edges encode communication and control flow—including conditional looping edges to implement iterative processes.

The core multi-agent feedback workflow (Evaluator-Optimizer loop) leverages LangGraph to allow repeated cycles of generation, validation, and refinement until convergence.

\section{Knowledge Base Construction}

\subsection{Web Scraping}

The accuracy and reliability of RAG-driven LLM workflows depend directly on the quality of the supporting knowledge base. The current work employs both automated and manual web scraping methods to continuously aggregate vendor-maintained documentation, whitepapers, and best practice guides from AWS, Azure, and potentially other providers, ensuring an up to date corpus.

\subsection{Docling: AI Document Parsing}

Complex vendor documents often appear in multi-format layouts (PDF, DOCX, PPTX), posing challenges for conventional parsers. Docling \cite{docling2025paper} is an AI-powered, open source document conversion tool that incorporates state-of-the-art models to parse such complex documents, extracting structured text, tables, and images accurately.

In the proposed framework, Docling integrates as a critical preprocessing step to generate clean, semantically meaningful document fragments suitable for downstream embedding and indexing \cite{docling2024web}.

\section{Local Model Development Environment}

\subsection{Ollama}

Ollama \cite{ollama2024web} streamlines the deployment of open-source LLMs such as Llama 3 and embedding-focused models like Nomic-Embed. The platform provides accessible APIs and abstracts hardware complexity, enabling efficient local embedding generation fundamental to RAG vector indices.

\subsection{ChromaDB: Vector Embedding Storage and Retrieval}

The framework utilizes the \textbf{ChromaDB} vector database \cite{chromadb2024} to store vector embeddings efficiently and facilitate fast similarity searches essential for RAG.

ChromaDB is a purpose-built, open-source vector database optimized for machine learning workloads, allowing scalable and low-latency retrieval of relevant documents based on approximate nearest neighbor searches. By indexing the Ollama-generated embeddings (derived from the \texttt{nomic-embed-text} model), ChromaDB enables Validator agents to retrieve contextually pertinent knowledge fragments rapidly, grounding assessments with high precision.

The integration of ChromaDB in the knowledge retrieval pipeline is critical for system performance and accuracy. The database supports complex queries where exact keyword matching fails, leveraging semantic similarity to connect diverse phrasings and domain-specific terminology prevalent in cloud service descriptions.

Thus, ChromaDB serves as the backbone vector store enabling reliable, validated factual grounding and iterative self-correction within the Evaluator-Optimizer architecture deployed.

\subsection{Cross-Encoder Reranking}

While vector databases excel at fast, approximate retrieval, pure bi-encoder semantic similarity often struggles to distinguish highly nuanced domain terminology. Gao et al.~\cite{gao2023retrieval} survey advanced RAG architectures and identify two-stage retrieval as a best practice: a fast embedding-based search fetches a broad candidate set, after which a computationally intensive cross-encoder rescores and reranks the candidates for contextual relevance. The cross-encoder attends jointly to the query and each candidate passage, capturing asymmetric relevance signals that bi-encoder scoring misses. In the implemented system, the \texttt{FlashRank} library (using \texttt{ms-marco-MiniLM-L-12-v2}) performs this reranking step, returning only the top-$k$ highest-scoring excerpts to validator agents. For cloud architecture validation---where a query about ``DynamoDB on-demand capacity'' must retrieve provisioning details rather than generic NoSQL introductions---this precision improvement is material.

\subsection{Apple Metal Performance Shaders (MPS)}

To achieve high throughput on commodity hardware without relying on costly GPUs, the system utilizes Apple’s Metal Performance Shaders (MPS) framework \cite{apple2024mps}. MPS allows ML workloads, particularly PyTorch-based embedding model computations, to run natively and efficiently on Apple Silicon GPUs, optimizing local development cycles.

\section{Review of Existing Work}

\subsection{Advanced RAG and Self-Correction}

Building on the original RAG framework \cite{lewis2020rag}, recent literature presents improvements emphasizing self-correction via reflection \cite{asai2023selfrag}. Such methods position self-evolving LLMs that critique and refine outputs, directly motivating the Evaluator-Optimizer approach to architecture recommendation in the thesis \cite{anthropic2024agents}.

\subsection{LLM-as-a-Judge Evaluation Frameworks}

Traditional evaluation metrics for natural language generation, such as ROUGE-L or METEOR, measure surface-level token overlap and struggle when multiple valid architectural solutions exist for the same requirement. Phase 1 results underscored this limitation: the agentic system produced more technically precise outputs than the plain LLM baseline, yet scored lower on ROUGE-L because it used service-specific vocabulary absent from the reference text.

Zheng et al.~\cite{zheng2024judging} formalise the LLM-as-a-judge paradigm, demonstrating that a strong frontier model can reliably score open-ended responses on expert criteria, achieving high agreement with human annotators when given carefully structured rubrics. Crucially, they identify that the choice of judge model matters: a judge drawn from the same model family as the system under test risks self-serving bias. This finding informed the evaluation design in Phase 2, where an independent judge model scores outputs across six technical dimensions (completeness, accuracy, security, scalability, best practices adherence, and specificity) on a 1--10 scale.

\subsection{Multi-Agent Frameworks for LLMs}

Frameworks such as AutoGen \cite{wu2023autogen} enable multi-agent conversational workflows using LLMs, providing flexible agent orchestration but typically without explicitly defined cyclic state management required for iterative validation loops. The proposed architecture advances such designs by employing LangGraph \cite{langgraph2024docs} for granular and stateful multi-agent control.

\subsection{AI in Cloud Environment Management}

Recent proprietary solutions including Amazon Q \cite{aws2023q} and Microsoft Copilot for Azure \cite{microsoft2024copilot} exemplify AI integration in cloud ecosystems. Despite the inherent capabilities, proprietary tools operate within single cloud ecosystems and lack comprehensive multi-cloud solutions or architectural debate mechanisms—a gap the present research aims to fill.

\section{Identifying the Research Gap}

The conducted literature synthesis uncovers critical unmet needs:

\begin{enumerate}
    \item \textbf{Lack of Comprehensive Architectural Tools:} Current AI assistants provide vendor-specific, fragmented guidance lacking end-to-end validated cloud solution synthesis \cite{aws2023q, microsoft2024copilot}.
    \item \textbf{Retrieval Quality in Technical Domains:} Standard vector search is insufficient for cloud documentation, where queries involve dense, overlapping service terminology. Two-stage retrieval with cross-encoder reranking~\cite{gao2023retrieval} is needed but absent from existing cloud AI tools.
    \item \textbf{Evaluation Methodology Mismatch:} Lexical metrics penalise technically correct but paraphrastically different outputs, a problem Zheng et al.~\cite{zheng2024judging} resolve through structured LLM judgment---yet this approach remains absent from cloud architecture evaluation practice.
    \item \textbf{Potential for Reliability via Advanced RAG and Self-Reflection:} Emerging methods \cite{lewis2020rag, asai2023selfrag} empower factual grounding and iterative self-correction, validating the Evaluator-Optimizer cyclical feedback methodology \cite{anthropic2024agents}.
\end{enumerate}

The present research targets these gaps by engineering an agentic multi-provider framework that combines cyclical validation, two-stage reranked retrieval, structured LLM judgment, and an extensible provider registry---creating a foundation that could ultimately accommodate any cloud vendor meeting the documented interface.